\section{Limitations}
% === ROUND 2 --- §6 Limitations: goal & contents ==========================
% GOAL: pre-empt the obvious objections, each tied to a result. CONTAINS: absolute unique
% R^2 is a few percent of the noise ceiling (a scale fact, not a weakness); "real beyond the
% nuisance ACTUALLY subtracted" (surprisal held out by design); the powered claim is one
% subject/cohort; the Q1 quality entanglement is unresolved at this cheap benchmark; Q2 is
% underpowered at the fold unit.
% ===========================================================================
The UTS03 result is conditional on one deeply sampled participant.
Its spatially dependent voxels provide resolution, not population replication, and its five story folds share training data \evd{E006}.
The nuisance set subtracts the listed rate, duration, frequency, and static-semantic variables; surprisal and imageability remain outside it, so “beyond nuisance” means beyond those variables actually controlled.

E004 has only five seed-averaged folds and uses an imperfect block-permutation sensitivity control.
Its corrected interval cannot exclude a small lever near the observed mean \evd{E004}.
E008 evaluates nine participants exposed to the same 1000 sentences, so the participant and stimulus axes are not jointly large \evd{E008}.
E011 changes adapter capacity at the same effective operating point rather than inducing a stronger representation change, and E013's voxelwise result is a failed MSE-distillation manipulation rather than a substrate-wide impossibility \evd{E011}\evd{E013}.

E009 measures OOD perplexity only, with per-arm language quality reported rather than exactly equalized; because its representation fulcrum is unreliable, it does not identify the payoff of a successful manipulation \evd{E009}.
E024 is a precondition result: no five-arm privileged-information intervention was completed, and its seven-paragraph corpus cannot support the intended stable paragraph-disjoint learning curve \evd{E024}.
